\documentclass[11pt]{article}
\usepackage[margin=1in]{geometry}
\usepackage{amsmath,amssymb,amsthm,graphicx,microtype}
\usepackage[colorlinks=true,linkcolor=blue,urlcolor=blue,citecolor=blue]{hyperref}
\newtheorem{theorem}{Theorem}
\newtheorem{remark}{Remark}
\newcommand{\R}{\mathbb R}
\newcommand{\T}{\mathfrak T}
\newcommand{\norm}[1]{\left\lVert #1\right\rVert}
\title{A Gaussian counterexample to the mixed-norm tensor identity\\in arXiv:2405.09378}
\author{Autonomous mathematical discovery run \texttt{fa\_banach\_001}}
\date{11 August 2026}

\begin{document}
\maketitle

\noindent\textbf{Status.}
\texttt{candidate\_counterexample\_likely\_valid}.  This packet disproves
the explicit identity used to motivate the source's final conjectural route;
it does not solve the broader, non-theorem-shaped program of characterizing
all mixed norms of metaplectic Wigner distributions.

\section{Source statement}

The source is Giacchi's arXiv:2405.09378, later published in the
\emph{Journal of Fourier Analysis and Applications} \cite{Giacchi2024}.
On page 28 of the arXiv PDF, equation (57) in the final remark states that,
for
\[
 M=\begin{pmatrix}
 I+P_{12}Q_{12}^{T}&-P_{12}\\
 -Q_{12}^{T}&I
 \end{pmatrix},
\]
the rescaling operator has the tensor-product property
\begin{equation}\label{source-identity}
 \norm{\T_M(f\otimes g)}_{L^{p,q}}
 =\norm f_p\norm g_q.
\end{equation}
The remark conjectures that this identity can be used to study and possibly
characterize mixed norms of metaplectic Wigner distributions.  The published
version contains the same assertion as equation (42) in Remark 8.5.

\begin{figure}[ht]
\centering
\includegraphics[width=\textwidth]{figures/open_problem_crop.png}
\caption{The full-width source excerpt, arXiv PDF page 28.}
\end{figure}

\section{Definitions and proof intuition}

The source uses
\[
 \norm F_{L^{p,q}}
 =\left(\int_{\R^d}\left(\int_{\R^d}|F(x,\xi)|^p\,dx\right)^{q/p}d\xi\right)^{1/q}
\]
for finite positive $p,q$, and
$\T_LF(z)=|\det L|^{1/2}F(Lz)$.  We need only finite exponents.

\paragraph{Proof intuition.}
Choose $P_{12}=0$.  Then $M$ is a determinant-one lower shear, so a tensor
$f(x)g(\xi)$ becomes $f(x)g(\xi-Q^Tx)$.  An ordinary $L^p$ norm cannot see
such a measure-preserving shear, but a mixed norm with unequal exponents does:
the inner Gaussian integration changes the covariance of the outer Gaussian.
The determinant of that covariance change is exactly
$\det(I+Q^TQ)$, producing a nontrivial factor whenever $p\ne q$.

\section{Counterexample and exact lower-shear calculation}

\begin{theorem}\label{thm:counterexample}
Let $d\ge1$, $0<p,q<\infty$, and $Q\in\R^{d\times d}$.  Set
\[
 P_{12}=0,
 \qquad
 M_Q=\begin{pmatrix}I&0\\-Q^T&I\end{pmatrix},
 \qquad
 \gamma(x)=e^{-\pi|x|^2}.
\]
Then, with $D=\det(I+Q^TQ)$,
\begin{equation}\label{exact-ratio}
 \frac{\norm{\T_{M_Q}(\gamma\otimes\gamma)}_{L^{p,q}}}
 {\norm\gamma_p\norm\gamma_q}
 =D^{\frac1{2q}-\frac1{2p}}.
\end{equation}
Consequently, if $Q\ne0$ and $p\ne q$, the source identity
\eqref{source-identity} is false.  In particular, $d=1$, $Q=1$, $p=1$,
and $q=2$ give the ratio $2^{-1/4}$.
\end{theorem}

\begin{proof}
The shear has determinant one, hence
\[
 \T_{M_Q}(\gamma\otimes\gamma)(x,\xi)
 =\exp\bigl(-\pi(|x|^2+|\xi-Q^Tx|^2)\bigr).
\]
Put $A=I+QQ^T$.  Completing squares gives
\begin{align*}
 |x|^2+|\xi-Q^Tx|^2
 &=(x-A^{-1}Q\xi)^TA(x-A^{-1}Q\xi)\\
 &\quad+\xi^T(I+Q^TQ)^{-1}\xi.
\end{align*}
Here we used the standard identities
$I-Q^TA^{-1}Q=(I+Q^TQ)^{-1}$ and
$\det A=\det(I+Q^TQ)=D$.  Gaussian integration in $x$ therefore yields
\[
 \norm{\T_{M_Q}(\gamma\otimes\gamma)(\cdot,\xi)}_p
 =p^{-d/(2p)}D^{-1/(2p)}
 \exp\bigl(-\pi\xi^T(I+Q^TQ)^{-1}\xi\bigr).
\]
Taking the outer $L^q$ norm gives
\[
 \norm{\T_{M_Q}(\gamma\otimes\gamma)}_{L^{p,q}}
 =p^{-d/(2p)}q^{-d/(2q)}D^{-1/(2p)+1/(2q)}.
\]
On the other hand,
$\norm\gamma_p\norm\gamma_q=p^{-d/(2p)}q^{-d/(2q)}$, proving
\eqref{exact-ratio}.  If $Q\ne0$, the positive semidefinite matrix $Q^TQ$
is nonzero, so $D>1$; unequal exponents then make the last power nonzero.
\end{proof}

\begin{remark}[Exact correction within the chosen subfamily]
For $P_{12}=0$, identity \eqref{source-identity} holds for all tensor inputs
when $p=q$, because $L^{p,p}=L^p$ and $M_Q$ is measure preserving.  It also
holds when $Q=0$, when $M_Q=I$.  Theorem \ref{thm:counterexample} proves that
for $0<p,q<\infty$, these are the only possibilities if the identity is
required for every tensor input.
\end{remark}

\section{Verification, scope, and novelty}

The proof is exact and uses no numerical approximation.  A reviewer should
check two convention-sensitive points: the source puts the $L^p$ norm in the
$x$ variable inside the $L^q$ norm in $\xi$, and the determinant factor in
$\T_M$ is one for $M_Q$.  Both match the source definitions.

A bounded search on 11 August 2026 used the exact equation phrase, the title,
arXiv id 2405.09378, ``mixed norm tensor products,'' and ``erratum.''  It found
the arXiv and final journal versions, both containing the assertion, but no
erratum or prior correction.  Novelty confidence is moderate because this is
an elementary Gaussian calculation that may be informally known.

The result does not characterize $\norm{W_{\mathcal A}(f,g)}_{L^{p,q}}$ in
general.  It shows that the proposed route cannot begin from
\eqref{source-identity} without additional hypotheses or a replacement
inequality.

\paragraph{Human-review recommendation.}
Accept as a likely-valid counterexample to equation (57)/(42), after checking
the mixed-norm order.  Do not count it as a full solution of the broad final
research program.

\begin{thebibliography}{9}
\bibitem{Giacchi2024}
G.~Giacchi,
\emph{Boundedness of metaplectic operators within $L^p$ spaces, applications
to pseudodifferential calculus, and time-frequency representations},
J. Fourier Anal. Appl. \textbf{30} (2024), Article 69;
arXiv:2405.09378, \url{https://arxiv.org/abs/2405.09378}.
\end{thebibliography}

\end{document}

